\chapter{Where to Place the Adapters}
\label{ch:placement}

Chapter~\ref{ch:lora_ft} placed adapters on every block. This chapter asks what happens when only half of the blocks
receive attention adapters (rank 8, learning rate $10^{-4}$): the first half, the second half, the even blocks or the
odd blocks. The four placements separate \emph{where} the adapters start (the first half and the even blocks start
at block 0, the odd blocks at block 1, the second half in the middle of the model) from \emph{how they are spread}.
Table~\ref{tab:place} gives accuracy and cost.

\begin{table}[!htb]
\centering\footnotesize
\setlength{\tabcolsep}{3pt}
\begin{tabular}{lcccccc}
\toprule
 & first & test acc. & vs all blocks & vs full FT & peak & epoch \\
adapted blocks & block & (\%) & (pp, paired) & (pp, paired) & (GiB) & (s) \\
\midrule
\rawinput{data/tab_place.tex}
\bottomrule
\end{tabular}
\caption{Attention LoRA (rank 8, learning rate $10^{-4}$) on different sets of blocks; mean $\pm$ standard deviation
over 5 seeds. First block = first adapted block. vs all blocks, vs full FT = paired difference to the run with the
same seed. Peak = peak training memory at batch size 64 (identical in all seeds); epoch = train epoch time. Per seed: Appendix~\ref{app:seeds}; forward and backward times:
Table~\ref{tab:time}.}
\label{tab:place}
\end{table}

\section{Accuracy}

\paragraph{Adapting only the second half costs accuracy on both models.} The second half is 1.56 $\pm$ 0.15~pp
below all blocks on ViT-B and 1.18 $\pm$ 0.19~pp below on Swin-B, in every seed. On both models it also falls below
full fine-tuning (by 1.77 and 0.70~pp, in every seed), and its validation loss is the highest of all placements. In
this setting adapters on the later half alone are not enough. This is the opposite of the usual practice when only
part of a pretrained model is trained, which is to train the last layers (in the extreme case only the classifier),
on the view that early layers compute general features and later layers features specific to the task; Lee et
al.~\cite{lee2023surgical} describe this as the common approach. A possible explanation, which we have not tested, is that CIFAR-100 images
upsampled from $32\times32$ differ from ImageNet images mainly in their fine structure, which the early blocks
process; this would agree with their observation that the early layers are the ones to train when it is the
input that changes. The two ``second halves'' are not the same thing structurally (blocks 12--23 of
Swin-B are the last ten blocks of its third stage and its fourth stage), and the size of the loss depends on the
model (Chapter~\ref{ch:budget}).

\paragraph{The best of the other placements differs between the models.} On ViT-B the first half is the most
accurate placement, 0.21 $\pm$ 0.06~pp above all blocks in every seed. On Swin-B it is 0.38 $\pm$ 0.20~pp
\emph{below} all blocks in every seed. Both differences are detectable. We have not tested the cause; the first
half of Swin-B has the fewest adapter parameters of the four placements and no adapters in the fourth stage.

\paragraph{Interleaved placements lose little or nothing against all blocks.} Neither the even nor the odd blocks
differ detectably from all blocks on either model. Even blocks: $-$0.06 $\pm$ 0.09~pp on ViT-B and $-$0.03 $\pm$
0.42~pp on Swin-B. Odd blocks: $-$0.22 $\pm$ 0.23~pp on ViT-B (lower in 4 of 5 seeds) and $+$0.08 $\pm$ 0.27~pp on
Swin-B. No detectable difference is not the same as no difference: for the odd blocks on ViT-B the mean is negative
and the validation loss is slightly worse in all five seeds, so a small loss cannot be excluded there. Against full
fine-tuning, the interleaved placements stay 0.45--0.56~pp ahead on Swin-B and are 0.26--0.42~pp behind on ViT-B.

All placements were trained at the same learning rate and rank; a placement with fewer or later adapters may prefer
a different rate, so these differences hold for this setting.

\section{Training memory and time}

\begin{figure}[!htb]
\centering
% generated by scripts/make_figures.py; do not edit
\begin{tikzpicture}
\begin{groupplot}[group style={group size=2 by 1, horizontal sep=\figSep, vertical sep=1.3cm}, width=\figW, height=\figH, repaxis, ymin=0, ymajorgrids, xlabel={first adapted block $k$}, legend style={font=\footnotesize, draw=none, fill=none, legend cell align=left}, legend pos=south west]
\nextgroupplot[title={ViT-B and DeiT-B (12 blocks)}, ylabel={peak training memory (GiB)}, xtick={0,2,4,6,8,10}]
\addplot[cAttn, thick, mark=*, mark size=1.8pt] coordinates {(0,5.459) (1,5.059) (2,4.660) (3,4.260) (4,3.860) (5,3.461) (6,3.061) (8,2.262) (10,1.462)};
\addlegendentry{DeiT-B (profiler)}
\addplot[cAccent, only marks, mark=square, mark size=3pt, thick] coordinates {(0,5.459) (3,4.260) (6,3.061)};
\addlegendentry{ViT-B (training runs)}
\nextgroupplot[title={Swin-B (24 blocks)}, xtick={0,2,4,8,12,16,22}]
\addplot[cAttn, thick, mark=*, mark size=1.8pt] coordinates {(0,9.549) (2,7.175) (4,6.000) (8,4.854) (12,3.707) (16,2.560) (22,1.815)};
\draw[cGrey, densely dotted] (axis cs:2,0) -- (axis cs:2,10);
\node[font=\scriptsize, cGrey, rotate=90, anchor=south west] at (axis cs:2,0.2) {stage 2};
\draw[cGrey, densely dotted] (axis cs:4,0) -- (axis cs:4,10);
\node[font=\scriptsize, cGrey, rotate=90, anchor=south west] at (axis cs:4,0.2) {stage 3};
\draw[cGrey, densely dotted] (axis cs:22,0) -- (axis cs:22,10);
\node[font=\scriptsize, cGrey, rotate=90, anchor=south west] at (axis cs:22,0.2) {stage 4};
\end{groupplot}
\end{tikzpicture}

\caption{Peak training memory at batch size 64 against the first adapted block $k$, for attention LoRA on every
block from $k$ to the last. Left: DeiT-B, which has the shape of ViT-B (profiler), and three ViT-B training runs.
Right: Swin-B; the dotted lines mark the first block of stages 2, 3 and 4.}
\label{fig:mem_start}
\end{figure}

\paragraph{Memory is set by the first adapted block.} On ViT-B, the first half and the even blocks have their
adapters in different places, but both start at block 0, and their peak memory is identical (5.233~GiB). The odd
blocks start at block 1 and need 0.362~GiB less (4.871~GiB). The second half starts at block 6 and needs exactly six
times that amount less (3.061~GiB). This is the behaviour described in Section~\ref{sec:bg_memory}: the backward pass
must reach the first adapter, every block from there on keeps its intermediate results, and the blocks before it
keep nothing. Figure~\ref{fig:mem_start} shows the same rule when the first adapted block is moved one block at a
time: the peak falls by the same amount at every step on ViT-B and DeiT-B, and fastest in the first blocks of
Swin-B, which work on its largest tensors (Table~\ref{tab:swin_stages}). On Swin-B, the even and the odd blocks
adapt one block of every pair in every stage, so for memory they differ only in where they start (they also differ
in the type of attention window, regular or shifted, which does not affect memory): starting at block 1 instead of
block 0 saves 0.98~GiB.

The odd blocks and the second half of ViT-B have the same number of trainable parameters (298,084) and differ by
1.81~GiB, so the parameter count does not predict training memory.

\paragraph{Additional adapted blocks are cheap.} Once the first adapted block is fixed, adapting more blocks after
it adds little, because those blocks store their intermediate results anyway. On ViT-B, adapting all twelve blocks
needs 0.226~GiB more than adapting the first six: about 0.04~GiB per additional block, a tenth of the 0.36~GiB that
starting one block earlier costs (0.40~GiB when the earlier block also receives an adapter, as in
Figure~\ref{fig:mem_start}).

\paragraph{The backward pass follows the same rule.} It takes about half as long when it stops in the middle of the
model (116 against 237~ms per batch on ViT-B, 117 against 271~ms on Swin-B), while the forward pass, which always
runs through the whole model, hardly changes. The shorter training epochs of the later-starting placements
(Table~\ref{tab:place}) therefore come from the backward pass.

\section{The trade-off}

\begin{figure}[!htb]
\centering
% generated by scripts/make_figures.py; do not edit
\begin{tikzpicture}
\begin{groupplot}[group style={group size=2 by 1, horizontal sep=\figSep, vertical sep=1.3cm}, width=\figW, height=\figH, repaxis, ymajorgrids, xmajorgrids, xlabel={peak training memory at batch 64 (GiB)}, legend style={font=\scriptsize, draw=none, fill=none, at={(-0.12,-0.3)}, anchor=north, legend columns=-1, /tikz/every even column/.append style={column sep=8pt}}]
\nextgroupplot[title={ViT-B}, xmin=2, xmax=8.8, ylabel={test accuracy (\%)}]
\addplot[cAttn, only marks, mark=*, mark size=2.2pt, forget plot, error bars/.cd, y dir=both, y explicit] coordinates {(5.459,91.102) +- (0,0.077)};
\node[font=\scriptsize, anchor=west, inner sep=3pt] at (axis cs:5.459,91.102) {all};
\addplot[cAttn, only marks, mark=square*, mark size=2.2pt, forget plot, error bars/.cd, y dir=both, y explicit] coordinates {(5.233,91.314) +- (0,0.109)};
\node[font=\scriptsize, anchor=south west, inner sep=3pt] at (axis cs:5.233,91.314) {first half};
\addplot[cAttn, only marks, mark=triangle*, mark size=2.2pt, forget plot, error bars/.cd, y dir=both, y explicit] coordinates {(5.233,91.046) +- (0,0.075)};
\node[font=\scriptsize, anchor=east, inner sep=3pt] at (axis cs:5.233,91.046) {even};
\addplot[cAttn, only marks, mark=diamond*, mark size=2.2pt, forget plot, error bars/.cd, y dir=both, y explicit] coordinates {(4.871,90.886) +- (0,0.215)};
\node[font=\scriptsize, anchor=north, inner sep=3pt] at (axis cs:4.871,90.886) {odd};
\addplot[cAttn, only marks, mark=pentagon*, mark size=2.2pt, forget plot, error bars/.cd, y dir=both, y explicit] coordinates {(3.061,89.544) +- (0,0.127)};
\node[font=\scriptsize, anchor=west, inner sep=3pt] at (axis cs:3.061,89.544) {second half};
\addplot[cFT, only marks, mark=square*, mark size=2.2pt, forget plot, error bars/.cd, y dir=both, y explicit] coordinates {(8.174,91.310) +- (0,0.146)};
\node[font=\scriptsize, anchor=south east, inner sep=4pt] at (axis cs:8.174,91.310) {full FT};
\nextgroupplot[title={Swin-B}, xmin=2, xmax=15.3]
\addplot[cAttn, only marks, mark=*, mark size=2.2pt, forget plot, error bars/.cd, y dir=both, y explicit] coordinates {(9.549,91.744) +- (0,0.314)};
\node[font=\scriptsize, anchor=west, inner sep=3pt] at (axis cs:9.549,91.744) {all};
\addplot[cAttn, only marks, mark=square*, mark size=2.2pt, forget plot, error bars/.cd, y dir=both, y explicit] coordinates {(9.284,91.362) +- (0,0.197)};
\node[font=\scriptsize, anchor=north west, inner sep=3pt] at (axis cs:9.284,91.362) {first half};
\addplot[cAttn, only marks, mark=triangle*, mark size=2.2pt, forget plot, error bars/.cd, y dir=both, y explicit] coordinates {(9.156,91.718) +- (0,0.186)};
\node[font=\scriptsize, anchor=north east, inner sep=3pt] at (axis cs:9.156,91.718) {even};
\addplot[cAttn, only marks, mark=diamond*, mark size=2.2pt, forget plot, error bars/.cd, y dir=both, y explicit] coordinates {(8.180,91.822) +- (0,0.291)};
\node[font=\scriptsize, anchor=east, inner sep=3pt] at (axis cs:8.180,91.822) {odd};
\addplot[cAttn, only marks, mark=pentagon*, mark size=2.2pt, forget plot, error bars/.cd, y dir=both, y explicit] coordinates {(3.707,90.566) +- (0,0.188)};
\node[font=\scriptsize, anchor=west, inner sep=3pt] at (axis cs:3.707,90.566) {second half};
\addplot[cFT, only marks, mark=square*, mark size=2.2pt, forget plot, error bars/.cd, y dir=both, y explicit] coordinates {(14.121,91.264) +- (0,0.232)};
\node[font=\scriptsize, anchor=south east, inner sep=4pt] at (axis cs:14.121,91.264) {full FT};
\addlegendimage{cAttn, only marks, mark=*, mark size=2.2pt}
\addlegendentry{attention LoRA on the named blocks}
\addlegendimage{cFT, only marks, mark=square*, mark size=2.2pt}
\addlegendentry{full fine-tuning}
\addlegendimage{cGrey, mark=|, mark size=3pt}
\addlegendentry{$\pm$ std, 5 seeds}
\end{groupplot}
\end{tikzpicture}

\caption{Test accuracy against peak training memory at batch size 64 for attention LoRA on all blocks and on the
four placements, and for full fine-tuning.}
\label{fig:placement}
\end{figure}

Figure~\ref{fig:placement} puts accuracy and memory together.
\begin{itemize}
  \item \textbf{Second half.} It saves the most memory: 44\% against all blocks on ViT-B and 61\% on Swin-B. But it
        costs 1.2--1.6~pp and ends below full fine-tuning on both models.
  \item \textbf{Odd blocks.} Leaving out block 0 saves 11\% of the memory of all-blocks LoRA on ViT-B and 14\% on
        Swin-B (40\% and 42\% against full fine-tuning) and shortens the training epoch by 8\%, with no detectable
        loss of accuracy against all blocks on Swin-B; on ViT-B a small loss cannot be excluded. Most of the saving
        (60\% and 70\% of it) comes from not reaching block 0, the rest from adapting half as many blocks.
  \item \textbf{Placements that start at block 0} need almost as much memory as all blocks. Which of them is the
        most accurate differs between the models: the first half on ViT-B, the interleaved blocks on Swin-B.
\end{itemize}
Within LoRA, the choice of blocks thus trades training memory against accuracy, and the memory side follows the
same rule on both models. How accuracy changes between a start at block 1 and a start in the middle of the model is
taken up in Chapter~\ref{ch:budget}.
